\chapter{Chapter05 — Majority Fourier One Norm}
%%% generated by the blueprint pipeline from FABL.Chapter05.MajorityFourierOneNorm
\section{Overview}

Declarations of the Lean module \texttt{FABL.Chapter05.MajorityFourierOneNorm}.

\section{Declarations}

\begin{theorem}[Fourier norm of a degree part]
\lean{FABL.fourierOneNorm_degreePart_eq_sum_level}
\label{FABL.fourierOneNorm_degreePart_eq_sum_level}
\leanok
\uses{FABL.degreePart, FABL.fourierCoeff, FABL.fourierCoeff_degreePart, FABL.fourierOneNorm}
\difficulty{3}
Let $n,k\in\mathbb N$, and let $f:\{-1,1\}^n\to\mathbb R$. The Fourier one norm of the
degree-$k$ part of $f$ equals $\sum_{S\subseteq\{1,\ldots,n\},\,|S|=k}|\widehat f(S)|$,
where $\widehat f(S)$ is the Fourier coefficient indexed by $S$.
\end{theorem}

\begin{theorem}[Fourier norm as a sum over degrees]
\lean{FABL.fourierOneNorm_eq_sum_degreeParts}
\label{FABL.fourierOneNorm_eq_sum_degreeParts}
\leanok
\uses{FABL.degreePart, FABL.fourierOneNorm, FABL.fourierOneNorm_degreePart_eq_sum_level}
\difficulty{4}
Let $n\in\mathbb N$ and let $f:\{-1,1\}^n\to\mathbb R$. Its Fourier one norm equals the
sum, over degrees $k=0,1,\ldots,n$, of the Fourier one norms of the degree-$k$ parts of
$f$.
\end{theorem}

\begin{theorem}[Odd-degree Fourier norm of majority]
\lean{FABL.fourierOneNorm_degreePart_majority_odd}
\proofsource{Analysis-of-Boolean-Functions-by-Ryan-ODonnell}{
  pdf-ffbe9c7b3306-p122-b004,
  pdf-ffbe9c7b3306-p123-b001,
  pdf-ffbe9c7b3306-p123-b002
}
\label{FABL.fourierOneNorm_degreePart_majority_odd}
\leanok
\uses{FABL.BooleanFunction.toReal, FABL.degreePart, FABL.fourierCoeff, FABL.fourierCoeff_majority_two_mul_add_one, FABL.fourierOneNorm, FABL.fourierOneNorm_degreePart_eq_sum_level, FABL.majority}
\difficulty{5}
Let $m,j\in\mathbb N$ with $j\le m$. For majority on the sign cube of dimension $2m+1$,
the Fourier one norm of its degree-$(2j+1)$ part is
$\binom{m}{j}\,\frac{1}{2j+1}\,\frac{2m+1}{2^{2m}}\binom{2m}{m}$.
\end{theorem}

\begin{theorem}[Even-degree Fourier norm of majority]
\lean{FABL.fourierOneNorm_degreePart_majority_odd_eq_zero_of_even}
\proofsource{Analysis-of-Boolean-Functions-by-Ryan-ODonnell}{
  pdf-ffbe9c7b3306-p122-b003,
  pdf-ffbe9c7b3306-p122-b004
}
\label{FABL.fourierOneNorm_degreePart_majority_odd_eq_zero_of_even}
\leanok
\uses{FABL.BooleanFunction.toReal, FABL.degreePart, FABL.fourierCoeff_majority_eq_zero_of_odd_arity_of_even_card, FABL.fourierOneNorm, FABL.fourierOneNorm_degreePart_eq_sum_level, FABL.majority}
\difficulty{4}
Let $m,k\in\mathbb N$ with $k$ even. The Fourier one norm of the degree-$k$ part of
majority on the sign cube of dimension $2m+1$ is zero.
\end{theorem}

\begin{theorem}[Even-zero sum reduces to odd terms]
\lean{FABL.sum_range_eq_sum_odd_of_even_eq_zero}
\label{FABL.sum_range_eq_sum_odd_of_even_eq_zero}
\leanok
\difficulty{4}
Let $a$ be a real sequence such that $a_{2j}=0$ for every nonnegative integer $j$, and
let $m$ be a nonnegative integer. Then
\[\sum_{k=0}^{2m+1}a_k=\sum_{j=0}^{m}a_{2j+1}.\]
\end{theorem}

\begin{definition}[Binomial half odd reciprocal expectation]
\lean{FABL.binomialHalfOddReciprocalExpectation}
\statementsource{Analysis-of-Boolean-Functions-by-Ryan-ODonnell}{pdf-ffbe9c7b3306-p137-b003}
\label{FABL.binomialHalfOddReciprocalExpectation}
\leanok
For each natural number $m$, let $J$ have the binomial distribution with parameters $m$
and $1/2$. The binomial half odd reciprocal expectation is $\mathbb{E}[1/(2J+1)]$.
\end{definition}

\begin{theorem}[Binomial reciprocal expectation formula]
\lean{FABL.binomialHalfOddReciprocalExpectation_eq_sum}
\label{FABL.binomialHalfOddReciprocalExpectation_eq_sum}
\leanok
\uses{FABL.binomialHalfOddReciprocalExpectation}
\difficulty{4}
For every natural number $m$, the expectation of $1/(2J+1)$ when $J$ has the binomial
distribution with $m$ trials and success probability $1/2$ equals
$\displaystyle\sum_{j=0}^{m}\binom{m}{j}2^{-m}\frac{1}{2j+1}$.
\end{theorem}

\begin{theorem}[Fourier norm of odd majority]
\lean{FABL.fourierOneNorm_majority_odd_eq_binomialExpectation}
\label{FABL.fourierOneNorm_majority_odd_eq_binomialExpectation}
\leanok
\uses{FABL.BooleanFunction.toReal, FABL.binomialHalfOddReciprocalExpectation, FABL.binomialHalfOddReciprocalExpectation_eq_sum, FABL.degreePart, FABL.fourierOneNorm, FABL.fourierOneNorm_degreePart_majority_odd, FABL.fourierOneNorm_degreePart_majority_odd_eq_zero_of_even, FABL.fourierOneNorm_eq_sum_degreeParts, FABL.majority, FABL.sum_range_eq_sum_odd_of_even_eq_zero}
\difficulty{4}
For every nonnegative integer $m$, the Fourier one-norm of the real-valued majority
function on $2m+1$ sign coordinates equals
$\mathbb{E}[1/(2J+1)]\,((2m+1)/2^m)\binom{2m}{m}$, where $J$ has the binomial
distribution with parameters $m$ and $1/2$.
\end{theorem}

\begin{theorem}[Binomial coefficient quotient identity]
\lean{FABL.cast_choose_div_succ}
\label{FABL.cast_choose_div_succ}
\leanok
\difficulty{3}
For natural numbers $n$ and $k$, the following identity holds:
\[(\binom{n}{k})/ (k+1) = \binom{n+1}{k+1}/(n+1).\]
\end{theorem}

\begin{theorem}[Binomial sum with reciprocal denominators]
\lean{FABL.sum_cast_choose_div_succ}
\label{FABL.sum_cast_choose_div_succ}
\leanok
\uses{FABL.cast_choose_div_succ}
\difficulty{4}
For every nonnegative integer $m$,
$\displaystyle\sum_{j=0}^{m}\frac{\binom{m}{j}}{j+1}=\frac{2^{m+1}-1}{m+1}$.
\end{theorem}

\begin{theorem}[Even reciprocal binomial sum]
\lean{FABL.binomialHalfEvenReciprocalSum_eq}
\label{FABL.binomialHalfEvenReciprocalSum_eq}
\leanok
\uses{FABL.sum_cast_choose_div_succ}
\difficulty{3}
For every natural number $m$,
$\displaystyle\sum_{j=0}^{m}\binom{m}{j}\frac{1}{2^m}\frac{1}{2j+2}=\frac{1-2^{-(m+1)}}{m+1}$.
\end{theorem}

\begin{theorem}[Binomial coefficient quotient identity with two shifts]
\lean{FABL.cast_choose_div_two_succs}
\label{FABL.cast_choose_div_two_succs}
\leanok
\uses{FABL.cast_choose_div_succ}
\difficulty{4}
For natural numbers $m$ and $j$, the following identity holds:
\[\binom{m}{j}/((j+1)(j+2))=\binom{m+2}{j+2}/((m+1)(m+2)).\]
\end{theorem}

\begin{theorem}[Shifted binomial sum bound]
\lean{FABL.sum_shifted_choose_le_pow}
\label{FABL.sum_shifted_choose_le_pow}
\leanok
\difficulty{3}
For every nonnegative integer $m$,
\[
\sum_{j=0}^{m}\binom{m+2}{j+2}\le 2^{m+2}.
\]
\end{theorem}

\begin{theorem}[Binomial reciprocal sum bound]
\lean{FABL.binomialHalfTwoSuccReciprocalSum_le}
\label{FABL.binomialHalfTwoSuccReciprocalSum_le}
\leanok
\uses{FABL.cast_choose_div_two_succs, FABL.sum_shifted_choose_le_pow}
\difficulty{4}
For every natural number $m$,
$\displaystyle\sum_{j=0}^{m}\binom{m}{j}2^{-m}\frac{1}{(j+1)(j+2)}\leq
\frac{4}{(m+1)(m+2)}$.
\end{theorem}

\begin{theorem}[Odd reciprocal binomial expectation bounds]
\lean{FABL.binomialHalfOddReciprocalExpectation_bounds}
\label{FABL.binomialHalfOddReciprocalExpectation_bounds}
\leanok
\uses{FABL.binomialHalfEvenReciprocalSum_eq, FABL.binomialHalfOddReciprocalExpectation, FABL.binomialHalfOddReciprocalExpectation_eq_sum, FABL.binomialHalfTwoSuccReciprocalSum_le}
\difficulty{5}
For every natural number $m$, let $J$ have the binomial distribution with $m$ trials and
success probability $1/2$. Then $\displaystyle\frac{1-2^{-(m+1)}}{m+1}\le \mathbb
E\!\left[\frac{1}{2J+1}\right]\le\frac{1-2^{-(m+1)}}{m+1}+\frac{4}{(m+1)(m+2)}$.
\end{theorem}

\begin{theorem}[Odd Reciprocal Binomial Expectation]
\lean{FABL.tendsto_binomialHalfOddReciprocalExpectation_mul_oddArity_div_two}
\label{FABL.tendsto_binomialHalfOddReciprocalExpectation_mul_oddArity_div_two}
\leanok
\uses{FABL.binomialHalfOddReciprocalExpectation, FABL.binomialHalfOddReciprocalExpectation_bounds}
\difficulty{5}
Let $J_m$ have the binomial distribution with parameters $m$ and $1/2$. Then
$\frac{2m+1}{2}\,\mathbb{E}\!\left[\frac{1}{2J_m+1}\right]\to1$ as $m\to\infty$.
\end{theorem}

\begin{theorem}[Odd-majority Fourier norm and influence]
\lean{FABL.fourierOneNorm_majority_odd_eq_expectation_mul_influence}
\label{FABL.fourierOneNorm_majority_odd_eq_expectation_mul_influence}
\leanok
\uses{FABL.BooleanFunction.toReal, FABL.binomialHalfOddReciprocalExpectation, FABL.fourierOneNorm, FABL.fourierOneNorm_majority_odd_eq_binomialExpectation, FABL.majority, FABL.oddMajorityInfluence}
\difficulty{3}
For every nonnegative integer $m$, the Fourier one-norm of the real-valued majority
function on $2m+1$ sign coordinates equals $\mathbb{E}[1/(2J+1)](2m+1)2^m I_m$, where
$J$ has the binomial distribution with parameters $m$ and $1/2$, and
$I_m=\binom{2m}{m}/2^{2m}$.
\end{theorem}

\begin{theorem}[Majority fourier one norm main eq book main]
\lean{FABL.majorityFourierOneNormMain_eq_bookMain}
\statementsource{Analysis-of-Boolean-Functions-by-Ryan-ODonnell}{pdf-ffbe9c7b3306-p137-b003}
\label{FABL.majorityFourierOneNormMain_eq_bookMain}
\leanok
\uses{FABL.oddMajorityInfluenceMain}
\difficulty{4}
% dataset-pipeline: stub -- prose pending build_informal_statements.py
Statement pending regeneration.
\end{theorem}

\begin{theorem}[Asymptotic Fourier norm of odd majority]
\lean{FABL.fourierOneNorm_majority_odd_isEquivalent}
\label{FABL.fourierOneNorm_majority_odd_isEquivalent}
\leanok
\uses{FABL.BooleanFunction.toReal, FABL.binomialHalfOddReciprocalExpectation, FABL.fourierOneNorm, FABL.fourierOneNorm_majority_odd_eq_expectation_mul_influence, FABL.majority, FABL.majorityFourierOneNormMain_eq_bookMain, FABL.oddMajorityInfluence, FABL.oddMajorityInfluenceMain, FABL.tendsto_binomialHalfOddReciprocalExpectation_mul_oddArity_div_two, FABL.tendsto_oddMajorityInfluence_div_main}
\difficulty{5}
As $m$ tends to infinity through the natural numbers, the Fourier one-norm of majority
on $2m+1$ sign coordinates is asymptotic to
$\frac{2}{\sqrt{\pi}}\frac{1}{\sqrt{2m+1}}2^{(2m+1)/2}$.
\end{theorem}
